\section{The order book in the settlement window}\label{sec:book}

The full-year results describe prices and trades. On the \DesignSessions{} sessions of the
order book sample the limit order book is reconstructed, which shows how the supply of
liquidity in the window changes on expiry days and whether the order flow bears the marks of
directional pressure. Each measure is computed per security and session and compared between
each monthly expiry and the control session in the same month, the trading day before it.
Table~\ref{tab:descriptives} reports the averages and Table~\ref{tab:contrasts} compares the
expiry-minus-control differences of the securities with derivatives with those of the placebo
group. The full set of paired tests is reported in Appendix~\ref{app:tests}.

\input{generated/tables/descriptives}

\subsection{Liquidity supply}

The window is more expensive to trade in on expiry days. The quoted spread widens from
\SpreadMeanBpsControl{} to \SpreadMeanBpsExpiry{} basis points, the displayed notional in the
book falls from \BookVisibleCrControl{} to \BookVisibleCrExpiry{} crore, and sweeping the
displayed book with an order of \ProbeSizeCr{} crore moves the price by
\MoveBpsSmallExpiry{} basis points against \MoveBpsSmallControl{} on control sessions. The share
of orders cancelled within one second of entry rises from \LifePhantomRateControl{}\% to
\LifePhantomRateExpiry{}\%. The widening of the spread and the rise in fleeting orders are
specific to securities with derivatives: relative to the placebo group, the spread widens by an
additional \DidMeanspreadbpsPlacebo{} basis points ($p = \DidMeanspreadbpsPlaceboP{}$) and the
share of orders cancelled within a second rises by an additional
\DidPhantomorderratePlacebo{} ($p = \DidPhantomorderratePlaceboP{}$). The
cost of sweeping the book rises by a similar amount in both sets of securities.

The rise in cost is modest relative to the level. Sweeping the book is several times more
expensive in the illiquid group, at \MoveBpsSmallIlliquid{} basis points, than in the liquid
group, at \MoveBpsSmallLiquid{}, on any session. The cost of moving the price is set far more by
the security than by the day.

\subsection{Pressure or congestion}

A window that is busier because more participants trade in both directions and a window in
which the price is pushed in one direction both show wider spreads and more cancellations.
They differ in three respects. Under sustained pressure, returns over successive seconds lean
the same way and the variance ratio rises; signed order flow is persistent; and a larger part
of the price impact of order flow is permanent. Under heavier two-sided trading the variance
ratio falls, flow is no more persistent, and more of the impact is transitory.

None of the three measures moves toward pressure. The variance ratio is
\VarianceRatioExpiry{} on expiry sessions against \VarianceRatioControl{} on controls, the
persistence of signed flow is unchanged, and the permanent share of price impact is
\ImpactPermanentShareExpiry{} against \ImpactPermanentShareControl{}. Against the placebo group
two of them move in the direction of congestion: relative to the placebo group, the change in the
variance ratio of securities with derivatives is $\DidVarianceratioPlacebo{}$
($p = \DidVarianceratioPlaceboP{}$), and the change in the permanent share of price impact
$\DidPermanentsharePlacebo{}$ ($p = \DidPermanentsharePlaceboP{}$). The price impact of a given flow over ten seconds rises,
from \ImpactShortBpsControl{} to \ImpactShortBpsExpiry{} basis points, but less of it lasts.
Figure~\ref{fig:pressure} shows the distributions.

\begin{figure}[htbp]
  \centering
  \includegraphics[width=\textwidth]{pressure_measures.pdf}
  \caption{The three measures that distinguish directional pressure from heavier two-sided
    trading, expiry against control sessions. Under sustained pressure the variance ratio and
    the persistence of flow rise and a larger share of impact is permanent.}
  \label{fig:pressure}
\end{figure}

\subsection{The settlement price and the close}

On the order book sample the settlement price ends further from the last traded price on
expiry sessions, \AbsReversalBpsExpiry{} basis points against \AbsReversalBpsControl{}, and
less volume is traded within a basis point of the running average. Both are consistent with a
more volatile window rather than with trading that tracks the average. The first is larger in
securities with derivatives than in the placebo group, by \DidAbsreversalbpsPlacebo{} basis
points ($p = \DidAbsreversalbpsPlaceboP{}$). On the full-year sample, which covers every
session and uses session fixed effects, the corresponding comparison of the final minutes with
the settlement price is not specific to these securities (Section~\ref{sec:window}).

\input{generated/tables/contrasts}

\subsection{Concealment}

The share of resting size that is not displayed is high on every session, at
\ConcealedDepthShareControl{}\% on control sessions, and does not change on expiry. Concealed
orders become somewhat larger relative to displayed ones, from \ConcealedSizeMultipleControl{}
to \ConcealedSizeMultipleExpiry{} times their size, and are replenished slightly less often per share
traded. Fewer, larger concealed orders worked more slowly through the window describe the
execution of large expiring positions rather than an attempt to hide an accumulation aimed at
the settlement price.
